\section{Algebra Uniwersalna.}
Wejdziemy teraz na nieco wyższy poziom abstrakcji.
Szybko jednak zobaczymy,
że znacząco ułatwi nam to posługiwanie się strukturami,
które zdefiniowaliśmy.

Na czym polega problem?
Poznane przez nas struktury algebraiczne wyglądają wszystkie bardzo podobnie.
Mamy zbiór, pewne stałe i operacje wykonywane na tym zbiorze.
Możemy się już domyślać,
	że będziemy mieli wiele bardzo podobnych twierdzeń i definicji dla każdej ze zdefiniowanych struktur.

Przykładowo,
	w algebrze występuje pojęcie \textbf{homomorfizmu}---funkcji $h: A \to B$ między dwoma strukturami
	tego samego typu (np. między dwoma grupami, dwoma monoidami itd.),
	które dodatkowo spełniają pewne warunki.
Zwykle jednym z nich jest (jeśli $(A, \star, \dots)$ i $(B, \bullet, \dots)$ są jakimiś strukturami)
\begin{equation}
	h(a\star b) = h(a) \bullet h(b).
\end{equation}
Ale definicje homomorfizmów mają też więcej wymagań.
W przypadku monoidów $(A, \star, e_A)$, $(B, \star, e_B)$ chcemy aby zachodziło też
\begin{equation}
	h(e_A) = e_B.
\end{equation}
W przypadku grup podobna własność musi zachodzić też dla odwracalności.

Najczęściej homomorfizmy definiuje się osobno dla każdej struktury.
W trakcie studiów spotkaliśmy się już z przykładem homomorfizmu---mowa o przekształceniach liniowych
	(przestrzenie liniowe też są strukturą algebraiczną!).
Nasuwa się jednak pytanie: czy nie dałoby się uogólnić tego,
	abyśmy mieli tylko jedną definicję?
Chcielibyśmy mieć coś takiej postaci:
\\\\
\noindent \textbf{Definicja}.
	\textbf{Homomorfizmem} struktur $(A, \dots)$ i $(B, \dots)$ nazywamy funkcję $h: A \to B$,
	spełniającą warunki \dots
\\\\
Ciężko jednak powiedzieć, co należałoby wstawić w miejsce kropek, aby sformalizować tę definicję.
Podobna sytuacja zajdzie,
    kiedy zechcemy zacząć definiować
	\textbf{podprzestrzenie liniowe}, \textbf{podgrupy}, \textbf{podciała}, \textbf{podmonoidy} itd.
Definicje tych obiektów byłyby do siebie bardzo podobne.
Jak jednak zawrzeć je wszystkie w jednej definicji?
Odpowiedzi na te pytania dostarcza nam dziedzina algebry zwana \textit{Algebrą Uniwersalną}.

\subsection{Podstawowe definicje.}

\begin{definition}
	Niech $A$ będzie zbiorem.
	Mówimy, że funkcja $f: A^n \to A$ ma \textbf{arność} $n$ (od końcówki \textit{-arny},
	a to z łacińskiej końcówki \textit{-aris},
	w analogii do bin\textit{arny}, tern\textit{arny}, itd.).
\end{definition}

Ciekawa i przydatna własność zachodzi, gdy w powyższej definicji $n = 0$.
Zbiór $A^0$ definiuje się jako zbiór jednoelementowy (co jest tym jednym elementem nie ma znaczenia).
Funkcja $f:A^0 \to A$ jest więc funkcją ze zbioru jednoelementowego w zbiór $A$.
Możemy więc na nią patrzeć jak na pojedynczą stałą (bo zbiór wartości też jest jednoelementowy).

\begin{definition}
	\textbf{Algebrą} nazywamy trójkę $(\Sigma, \alpha, E)$,
	gdzie
	\begin{itemize}
		\item $\Sigma$ jest zbiorem, zwanym zbiorem \textbf{symboli},
		\item $\alpha: \Sigma \to \mathbb{N}$ jest funkcją, która każdemu symbolowi przypisuje arność,
		\item $E$ jest zbiorem równań używających zmiennych i symboli z $\Sigma$ jako funkcji.
	\end{itemize}
	Parę $(\Sigma, \alpha)$ nazywamy \textbf{sygnaturą} algebry $(\Sigma, \alpha, E)$.
\end{definition}

Definicja ta nie jest w pełni formalna (czym są ,,symbole"? czym są ,,zmienne"?
	co znaczy, że elementów $\Sigma$ używamy jako funkcji?
	jak możemy czemuś, co nie jest funkcją, przypisać arność?).
Sformalizowanie jej jednak wprowadziłoby więcej zamieszania niż pożytku,
	gdyż wymagałoby wprowadzenia mało istotnych, technicznych szczegółów.
Zamiast tego przejdziemy do rozjaśniającego przykładu.

\begin{example}
	Trójka $(\Sigma, \alpha, E)$, gdzie
	\begin{itemize}
		\item $\Sigma = \{f, g, e\}$,
		\item $\alpha: f \mapsto 2, g \mapsto 2, e \mapsto 0$,
		\item $E = \{g(x, y) = g(y, x), f(e, x) = x, f(x, e) = x\}$
	\end{itemize}
	jest algebrą.
\end{example}
Jak możemy interpretować, czym jest taka algebra?
Jest to nic innego jak definicja pewnej struktury algebraicznej,
gdzie mamy dwie operacje dwuargumentowe $f, g$ i jedną wyróżnioną stałą $e$,
że $g$ jest przemienne, a $e$ jest elementem neutralnym dla $f$.

Tłumacząc powyższy zapis na tradycyjną definicję, moglibyśmy napisać następująco.
\begin{definition}
	\textbf{Gałganem} nazywamy czwórkę $(A, f, g, e)$,
		gdzie $A$ jest zbiorem,
		$f, g: A^2 \to A$ są funkcjami dwuargumentowymi,
		$e \in A$ (lub, co w pewnym sensie jest równoważne, $e: \{0\} \to A$)
		i spełnione są następujące warunki:
	\begin{itemize}
		\item $\forall_{x\in A}\; f(e, x) = x,$
		\item $\forall_{x\in A}\; f(x, e) = x,$
		\item $\forall_{x,y\in A}\; g(x, y) = g(y, x).$
	\end{itemize}
\end{definition}

Zwykle zapis będziemy jednak upraszczać,
	podając jednocześnie zbiór symboli $\Sigma$ i funkcję $\alpha$ za jednym razem,
	n.p. w postaci
\begin{equation}
	\Sigma = \{f^2, g^2, e\},
\end{equation}
t.j. podając w definicji symboli ich arność w indeksie górnym, i pomijając ją dla $0$.
Idąc tą konwencją będziemy też zamiast $(\Sigma, \alpha, E)$ pisać $(\Sigma, E)$.

Algebry odpowiadają więc definicjom struktur.
Co odpowiada jednak obiektom, spełniającym te definicje?

\begin{definition}
    \textbf{Modelem algebry} $(\Sigma, \alpha, E)$ nazywamy zbiór $A$ (zwany \textbf{nośnikiem}),
	wraz z interpretacją symboli z $\Sigma$ jako funkcje o arnościach definiowanych przez $\alpha$ i spełniających równania $E$.
\end{definition}
Znowu nie jest to całkowicie formalna definicja,
na nasze potrzeby jednak w pełni wystarczy.

\begin{example}
	Zbiór $\mathbb{R}$ z działaniami $f(x, y) = x + y$, $g(x, y) = xy$ jest modelem gaugana.
\end{example}

Możemy teraz definiować nasze struktury w taki sam sposób jak gaugana.
\begin{definition}
	\textbf{Monoidem} nazywamy dowolny model algebry
	\begin{equation*}
		(\{\star^2, e\}, \{ x \star (y \star z) = (x \star y) \star z, x \star e = x, e \star x = x \}).
	\end{equation*}
\end{definition}

Ważne, aby rozróżniać teraz dwa pojęcia: algebrą jest \textit{definicja} monoidu.
Modelem algebry jest natomiast dowolny monoid,
	czyli obiekt,
	który spełnia tę definicję.

Algebra uniwersalna nie jest jednak taka uniwersalna,
	jak byśmy chcieli.
Jedna z podanych przez nas w poprzedniej sekcji definicji nie stanowi algebry.
Mowa tutaj o ciałach.
Problemem jest to,
	że funkcja $\mathrm{inv}$ nie musi być określona dla elementu neutralnego dodawania w ciele.
Jest to warunek,
	którego nie da się sformułować używając pojęcia algebry.
Z tego powodu ciała będziemy musieli traktować osobno od pozostałych struktur.

\subsection{Zadania.}

\begin{exercise}
	Zapisz definicję magmy w postaci algebry.
\end{exercise}
\begin{exercise}
	Zapisz definicję półgrupy w postaci algebry.
\end{exercise}
\begin{exercise}
	Zapisz definicję grupy w postaci algebry.
\end{exercise}
\begin{exercise}
	Zapisz definicję grupy abelowej w postaci algebry.
\end{exercise}
\begin{exercise}
	Zapisz definicję pierścienia w postaci algebry.
\end{exercise}

